% Generated draft: requires mathematical and source review.
\documentclass[11pt]{article}
\usepackage[margin=25mm]{geometry}
\usepackage{fontspec}
\setmainfont{Latin Modern Roman}
\usepackage{amsmath,amssymb,mathtools,unicode-math}
\setmathfont{Latin Modern Math}
\usepackage{enumitem,needspace,xurl,hyperref}
\setlength{\parindent}{0pt}
\setlength{\parskip}{5pt}
\setlength{\emergencystretch}{3em}
\newcommand{\N}{\mathbb{N}}
\newcommand{\Z}{\mathbb{Z}}
\newcommand{\Q}{\mathbb{Q}}
\newcommand{\R}{\mathbb{R}}
\newcommand{\C}{\mathbb{C}}
\newcommand{\F}{\mathbb{F}}
\newcommand{\A}{\mathbb{A}}
\newcommand{\PP}{\mathbb{P}}
\newcommand{\E}{\mathbb{E}}
\newcommand{\eps}{\varepsilon}
\newcommand{\dd}{\,\mathrm{d}}
\begin{document}
\section*{UnsolvedMath ID: 801087}
\textbf{Source problem code:} catalogue entry
\subsection{Definitions and mathematical statement}
Let \(d\) be a positive integer and let \(f=(f_1,\ldots,f_d):\C^d\to\C^d\) be a polynomial automorphism, meaning that both \(f\) and its inverse are polynomial maps. Write \(f^n\) for its \(n\)-fold iterate and define \(\deg(f)=\max_{1\le i\le d}\deg(f_i)\), using total polynomial degree. Under the usual interpretation of the symbol in the source, the first dynamical degree is
\[
\lambda(f)=\lim_{n\to\infty}\bigl(\deg(f^n)\bigr)^{1/n}.
\]
The supplied record does not itself define \(\lambda\), so this interpretation requires confirmation against the source.

For a nonnegative integer \(k\), the notation \(\deg(f^n)\asymp n^k\) means that there exist constants \(c,C>0\) and an integer \(n_0\ge1\), allowed to depend on \(f\), such that
\[
c n^k\le\deg(f^n)\le C n^k\qquad\text{for every integer }n\ge n_0.
\]
With the convention \(\N=\{0,1,2,\ldots\}\), the first question is whether
\[
\forall d\ge1\;\forall f\in\operatorname{Aut}_{\mathrm{poly}}(\C^d),\qquad
\lambda(f)=1\ \Longrightarrow\ \exists k\in\Z_{\ge0}\;\bigl(\deg(f^n)\asymp n^k\bigr).
\]
Here \(\operatorname{Aut}_{\mathrm{poly}}(\C^d)\) denotes the set of polynomial automorphisms just defined. The source writes \(k\in\N\) without specifying whether \(0\) is included; if it means positive integers, the quantified exponent set must instead be \(\Z_{\ge1}\). This convention remains unresolved.

The further question is: whenever \(\lambda(f)=1\) and \(\deg(f^n)\asymp n^k\) for an integer \(k\ge1\), must \(f\) preserve a rational fibration? The supplied record does not define this last phrase. A possible interpretation of preservation is the existence of an algebraic variety \(B\), a dominant rational map \(\pi:\C^d\dashrightarrow B\) with \(0<\dim B<d\), and a birational self-map \(g:B\dashrightarrow B\) satisfying
\[
\pi\circ f=g\circ\pi
\]
as rational maps. This is only a candidate interpretation: the record does not specify the dimension or rationality of \(B\), connectedness or rationality conditions on the general fiber, or whether preservation instead requires \(\pi\circ f=\pi\). Consequently, the fibration clause cannot be made unambiguous from the supplied record alone.
\subsection{Short English statement}
Question 39 asks whether every polynomial automorphism of \(\C^d\) with dynamical degree one has degree growth comparable to an integer power of the iteration number, and whether positive-power growth forces preservation of a rational fibration. The meanings of \(\lambda\), \(\N\), and rational-fibration preservation require source confirmation. The registry labels the question open; this draft does not independently verify that status.
\subsection{Sources}
\begin{itemize}
\item Local UnsolvedMath registry, permanent ID \(801087\), exact source problem code \texttt{catalogue entry}, title ``Henon Maps: Open Problems --- Question 39.'' The supplied record attributes the problem to Berger and Bedford, with proposed year 2023, and labels it open.
\item Berger and Bedford, \emph{Henon Maps: A List of Open Problems} (2023), Question 39. Supplied source URL: \url{https://arxiv.org/abs/2312.03907}. The supplied provenance records access on 2026-10-16 and extraction by \texttt{arxiv\_latex}; these are supplied provenance claims, not independently verified access or source checks.
\item Supplied upstream revision: \texttt{d371ee917ac9647fbb06cd6fc529deec6f3dda82}. The registry record was created and updated at \texttt{2026-10-16T14:00:00Z}.
\end{itemize}
\end{document}
